\documentclass[11pt]{letter}
\usepackage[utf8]{inputenc}
\usepackage{geometry}
\geometry{a4paper, margin=1in}

% --- Author Information ---
\address{
    Dr Jaewoo Joo \\
    School of Mathematics and Physics \\
    University of Portsmouth \\
    PO1 3FX, United Kingdom \\
    Email: jaewoo.joo@port.ac.uk
}

\begin{document}

\begin{letter}{
    To the Editors \\
    The Journal of Physical Chemistry Letters \\
    American Chemical Society
}

\opening{Dear Editors,}

I am pleased to submit our manuscript entitled \textbf{``Non-Equilibrium Dynamics of the Time-Dependent Excitonic Coupling in Fluorescent Protein Dimers''} for consideration as a Letter in \textit{the Journal of Physical Chemistry Letters}.

This work addresses a long-standing theoretical tension in biophysics: the observation of robust excitonic delocalization (Davydov splitting) in fluorescent protein dimers despite the highly decoherent, warm, and wet biological environment. Using a novel quantum-classical hybrid workflow, we quantify the excitonic coupling in the dimeric Venus homodimer and provide a physical mechanism that reconciles these seemingly contradictory observations.

Our work offers three primary contributions of high interest to the physical chemistry \& biology community such as
\begin{enumerate}
    \item \textbf{Near-Field Multipolar Effects:} We demonstrate that the standard point-dipole approximation underestimates the coupling strength by a factor of 5.6. By employing a full 3D transition-density coupling formalism, we find a robust coupling of $J = 74.38 \text{ cm}^{-1}$, aligning our theoretical results with recent experimental circular dichroism measurements.
    \item \textbf{Separation of Timescales:} We propose that excitonic coupling is initially protected by a non-equilibrium dielectric response. Because absorption occurs on a sub-femtosecond timescale, the coupling operates in the optical dielectric limit ($\epsilon \approx 1.78$) before the onset of bulk solvent relaxation ($\epsilon \approx 78$).
    \item \textbf{Stochastic Dynamics:} Using stochastic Schrödinger equation simulations, we visualize the transition from a delocalized exciton superposition to incoherent hopping, providing a clear temporal map of how spectroscopic signatures are imprinted before environmental dephasing dominates.
\end{enumerate}

Given the journal's focus on reporting significant scientific advances in physical chemistry that require rapid dissemination, we believe our findings will be of great interest to your readers. Thank you for your time and consideration of our work.
\\

{Sincerely,}\\
{Robson Christie, Cerys Murray, Youngchan Kim, Jaewoo Joo}
\end{letter}
\end{document}
